% Unit 102 English translation of chapter7.tex lines 2472--2678.
% Source: Methods of Algebra, Volume 2, commit
% 9a5803ff77dd3257484cb177f851a73770a59dd3 (CC BY 4.0).
% stable-unit-id: o014.aljabr2.chapter7.galois-descent-h1
% Coverage: complete Section 7.11, the relation between Galois descent and
% nonabelian H^1, ending immediately before the Exercises.

% segment-id: o014.aljabr2.chapter7.galois-descent-h1.h001
\section{The Relation between Galois Descent and
\texorpdfstring{$\Hm^1$}{H1}}\label{sec:Galois-H1}
\hypertarget{o014.aljabr2.chapter7.galois-descent-h1}{ }
% segment-id: o014.aljabr2.chapter7.galois-descent-h1.p001
Descent is a powerful method for classifying various algebraic or geometric
objects; its general formulation requires the language of algebraic geometry.
For ease of exposition, we make two simplifying assumptions in this section.
% segment-id: o014.aljabr2.chapter7.galois-descent-h1.l001
\begin{enumerate}
	\item We use only Galois descent
	(Theorem~\ref{prop:Galois-descent} and the discussion following it), rather
	than flat descent or more general techniques.
	\item As far as possible, this book uses algebraic language, especially that
	of linear algebra, to explain the examples. We must therefore narrow the
	scope of the classification problem. In fact, we shall consider only a
	certain class of tensors in finite-dimensional vector spaces.
\end{enumerate}

% segment-id: o014.aljabr2.chapter7.galois-descent-h1.p002
This viewpoint naturally leads to the nonabelian cohomology $\Hm^1$ of the
Galois group. The discussion below draws on the tools of
\S\ref{sec:nonabelian-coh} and \S\ref{sec:descent-Galois}. Since we need to
allow infinite Galois extensions, we shall also continue to use the language
of profinite groups; see \S\ref{sec:profinite-groups}. We first introduce the
necessary notation.

% segment-id: o014.aljabr2.chapter7.galois-descent-h1.cv001
\begin{convention}
	Throughout this section, fix a field $K$ and write
	$\otimes:=\otimes_K$. For a finite-dimensional $K$-vector space $V$, write
	$V^\vee$ for its dual space. For each $(p,q)\in\Z_{\geq0}^2$, introduce the
	notation
	% segment-id: o014.aljabr2.chapter7.galois-descent-h1.q001
	\[
	T^{p,q}V:=V^{\otimes p}\otimes(V^\vee)^{\otimes q}.
	\]
	Following the convention of differential geometry, the elements of
	$T^{p,q}V$ are called \emph{tensors of type $(p,q)$ on $V$}.
	\index{tensor}
\end{convention}

% segment-id: o014.aljabr2.chapter7.galois-descent-h1.p003
Every isomorphism $g:V\to V'$ between finite-dimensional $K$-vector spaces
naturally induces an isomorphism
$T^{p,q}V\rightiso T^{p,q}V'$; for convenience, this induced isomorphism is
still denoted by $g$. Moreover, finite dimensionality gives a canonical
isomorphism
$(V_1\otimes V_2)^\vee\simeq V_1^\vee\otimes V_2^\vee$.

% segment-id: o014.aljabr2.chapter7.galois-descent-h1.p004
In the classification problems considered in this section, the principal
objects are data of the following form, defined over the field $K$:
% segment-id: o014.aljabr2.chapter7.galois-descent-h1.q002
\begin{equation}\label{eqn:tensor-data}
	\mathbf{t}=(V,I,(t_i,p_i,q_i)_{i\in I}),
\end{equation}
where
% segment-id: o014.aljabr2.chapter7.galois-descent-h1.l002
\begin{itemize}
	\item $V$ is a finite-dimensional $K$-vector space;
	\item $I$ is a set (which may be empty);
	\item $(p_i,q_i)$ is a family of pairs of nonnegative integers, indexed by
	$i\in I$;
	\item $t_i\in T^{p_i,q_i}V$ is a family of tensors, indexed by $i\in I$.
\end{itemize}

% segment-id: o014.aljabr2.chapter7.galois-descent-h1.p005
In other words, we study a family of tensors of specified types in a
finite-dimensional vector space. Given another datum
% segment-id: o014.aljabr2.chapter7.galois-descent-h1.q003
\[
\mathbf{t}'=(V',I,(t'_i,p_i,q_i)_{i\in I}),
\]
if there is an isomorphism of $K$-vector spaces $g:V\rightiso V'$ such that
$gt_i=t'_i$ for every $i$, then $g$ is called an isomorphism between the two
data, and we write
% segment-id: o014.aljabr2.chapter7.galois-descent-h1.q004
\[
g\mathbf{t}=\mathbf{t}'\quad\text{or}\quad
g:\mathbf{t}\rightiso\mathbf{t}'.
\]

% segment-id: o014.aljabr2.chapter7.galois-descent-h1.p006
We wish to classify such data $\mathbf{t}$ up to isomorphism, or more
specifically to classify those $\mathbf{t}$ satisfying additional
conditions. These conditions must be algebraically expressible and preserved
by isomorphisms. Although the theoretical framework is confined to linear
algebra, its range of applications is already quite broad. Here are several
special cases.
% segment-id: o014.aljabr2.chapter7.galois-descent-h1.l003
\begin{enumerate}
	\item Consider data $(V,b)$, where
	$b\in(V^\vee)^{\otimes2}$ is a tensor of type $(0,2)$. Classifying these
	data is equivalent to classifying bilinear forms on $V$. We may also impose
	conditions such as symmetry, antisymmetry, or nondegeneracy on $b$; all of
	these are ``algebraic'' conditions. Multilinear forms may be treated
	similarly.
	\item We may also classify $(V,\varphi_1,\ldots,\varphi_n)$, where
	$\varphi_i\in V\otimes V^\vee$ is a tensor of type $(1,1)$. Since
	$V\otimes V^\vee\simeq\End_K(V)$, this amounts to classifying $V$ together
	with a family of linear endomorphisms. Various algebraic relations may be
	imposed, such as commutativity among the endomorphisms.
	\item Next consider data $(V,\mu)$, where
	$\mu\in V\otimes(V^\vee)^{\otimes2}$ is a tensor of type $(1,2)$, which may
	also be regarded as an element of $\Hom_K(V\otimes V,V)$. This amounts to
	classifying $V$ together with a bilinear ``multiplication''
	$V\times V\to V$. If the operation is required to be associative and to
	have a unit (the latter may be regarded as a tensor of type $(1,0)$), these
	conditions make $V$ into a $K$-algebra. Other conditions may be imposed on
	$\mu$ in order to study structures such as Lie algebras and Jordan
	algebras.
\end{enumerate}

% segment-id: o014.aljabr2.chapter7.galois-descent-h1.p007
More generally, \eqref{eqn:tensor-data} can be extended in a natural way to
treat the classification of a family of vector subspaces, or at least lines,
in $T^{p,q}V$. We shall not pursue this because of space constraints.

% segment-id: o014.aljabr2.chapter7.galois-descent-h1.d001
\begin{definition}
	For a finite-dimensional $K$-vector space $V$, write $\GL(V)$ for its group
	of linear automorphisms. For data $\mathbf{t}$ as in
	\eqref{eqn:tensor-data}, define the subgroup of $\GL(V)$
	% segment-id: o014.aljabr2.chapter7.galois-descent-h1.q005
	\[
	G(\mathbf{t}):=\{g\in\GL(V):g\mathbf{t}=\mathbf{t}\}.
	\]
\end{definition}

% segment-id: o014.aljabr2.chapter7.galois-descent-h1.p008
Choose a basis for $V$ and write the conditions $gt_i=t_i$ in coordinates
for every $i\in I$. Thus $G(\mathbf{t})$ is cut out inside $\GL(V)$ by a
family of polynomial equations. In the examples above, if $\mathbf{t}$
corresponds to a nondegenerate symmetric (respectively, antisymmetric)
bilinear form $b:V\times V\to K$, then the corresponding $G(\mathbf{t})$ is
the orthogonal group $\mathrm{O}(V,b)$ (respectively, the symplectic group
$\mathrm{Sp}(V,b)$). If $\mathbf{t}$ is the tensor of type $(0,\dim V)$ given
by the determinant, then the corresponding group is the special linear
group $\SL(V):=\{g:\det g=1\}$.

% segment-id: o014.aljabr2.chapter7.galois-descent-h1.d002
\begin{definition}\label{def:tL}
	For every field extension $L|K$ and $K$-vector space $V$, write
	$V_L:=L\dotimes{K}V$, regarded as an $L$-vector space. Through the canonical
	isomorphism of $L$-vector spaces
	% segment-id: o014.aljabr2.chapter7.galois-descent-h1.q006
	\[
	(T^{p,q}V)_L\simeq T^{p,q}(V_L),
	\]
	every datum $\mathbf{t}$ on $V$ of the form \eqref{eqn:tensor-data}
	induces a datum $\mathbf{t}_L$ on $V_L$, defined over the extension field
	$L$.
\end{definition}

% segment-id: o014.aljabr2.chapter7.galois-descent-h1.p009
Field extension does not change the numerical part of the data,
$(\dim V,I,(p_i,q_i)_i)$.

% segment-id: o014.aljabr2.chapter7.galois-descent-h1.p010
Recall that when $L|K$ is a Galois extension, $\Gal(L|K)$ acts smoothly and
semilinearly on $V_L$ (Lemma~\ref{prop:Galois-descent-action}); we write the
action on the left. In fact, it acts smoothly and semilinearly on every
$T^{p,q}(V_L)$ (Definition~\ref{def:semilinear-action}), ``diagonally'' on
tensor products and by $\check v\mapsto\check v\sigma^{-1}$ on the dual space
$V^\vee$ (the contragredient action, or ``inverse step''). Relative to the
isomorphism above, its invariant subspace is
% segment-id: o014.aljabr2.chapter7.galois-descent-h1.q007
\[
T^{p,q}(V_L)^{\Gal(L|K)}=T^{p,q}V.
\]
This is precisely Galois descent at the level of objects. At the level of
morphisms, it is expressed by the following result.

% segment-id: o014.aljabr2.chapter7.galois-descent-h1.lm001
\begin{lemma}
	Let $V$ and $W$ be finite-dimensional $K$-vector spaces. For every
	$h\in\Hom_L(V_L,W_L)$ and $\sigma\in\Gal(L|K)$, define
	% segment-id: o014.aljabr2.chapter7.galois-descent-h1.q008
	\[
	{}^\sigma h:=\sigma h\sigma^{-1}:V_L\to W_L.
	\]
	This map is again $L$-linear. The formula gives a smooth semilinear action
	of $\Gal(L|K)$ on the $L$-vector space $\Hom_L(V_L,W_L)$. This action is
	compatible with composition of linear maps, and
	% segment-id: o014.aljabr2.chapter7.galois-descent-h1.q009
	\[
	\Hom_L(V_L,W_L)^{\Gal(L|K)}=\Hom_K(V,W).
	\]
	Here $\Hom_K(V,W)$ is embedded in $\Hom_L(V_L,W_L)$ by
	$f\mapsto f_L:=\identity_L\otimes f$.
\end{lemma}
% segment-id: o014.aljabr2.chapter7.galois-descent-h1.pf001
\begin{proof}
	The reader may check directly that ${}^\sigma h$ is indeed $L$-linear.
	For smoothness, choose bases of $V$ and $W$ and identify the elements of
	$\Hom_L(V_L,W_L)$ with matrices. Then $\Gal(L|K)$ acts in the standard way
	on every matrix entry. Since $h$ involves only finitely many matrix entries,
	they are all contained in some finite subextension $E|K$ of $L|K$. Hence the
	open subgroup $\Gal(L|E)$ fixes $h$. The remaining assertions about the
	action $(\sigma,h)\mapsto{}^\sigma h$ are readily verified.

	The equality
	$\Hom_L(V_L,W_L)^{\Gal(L|K)}=\Hom_K(V,W)$ follows directly from Galois
	descent and is also easy to see directly from matrices.
\end{proof}

% segment-id: o014.aljabr2.chapter7.galois-descent-h1.p011
Taking the special case $V=W$, we see that $\Gal(L|K)$ acts on $\GL(V_L)$
by $(\sigma,g)\mapsto{}^\sigma g$, and that
$\GL(V_L)^{\Gal(L|K)}=\GL(V)$.

% segment-id: o014.aljabr2.chapter7.galois-descent-h1.lm002
\begin{lemma}
	For data $\mathbf{t}$ of the form \eqref{eqn:tensor-data}, the smooth action
	of $\Gal(L|K)$ on $\GL(V_L)$ preserves $G(\mathbf{t}_L)$. Thus
	$G(\mathbf{t}_L)$ becomes a smooth $\Gal(L|K)$-group in the sense of
	Remark~\ref{rem:pro-finite-coh-nonabelian}; its invariant subgroup satisfies
	% segment-id: o014.aljabr2.chapter7.galois-descent-h1.q010
	\[
	G(\mathbf{t}_L)^{\Gal(L|K)}=G(\mathbf{t}).
	\]
\end{lemma}
% segment-id: o014.aljabr2.chapter7.galois-descent-h1.pf002
\begin{proof}
	Let $g\in G(\mathbf{t}_L)$ and $\sigma\in\Gal(L|K)$. Then
	% segment-id: o014.aljabr2.chapter7.galois-descent-h1.q011
	\[
	({}^\sigma g)\mathbf{t}_L
	=({}^\sigma g)(\sigma\mathbf{t}_L)
	=\sigma g\sigma^{-1}(\sigma\mathbf{t}_L)
	=\sigma(g\mathbf{t}_L)
	=\sigma(\mathbf{t}_L)
	=\mathbf{t}_L.
	\]
	Thus the action of $\Gal(L|K)$ preserves $G(\mathbf{t}_L)$. The remaining
	assertions follow immediately.
\end{proof}

% segment-id: o014.aljabr2.chapter7.galois-descent-h1.p012
In practice, the classification of data \eqref{eqn:tensor-data} usually
becomes simpler after extending scalars to a sufficiently large field $L$.
For example, the isomorphism class of a nondegenerate quadratic form over an
algebraically closed field is completely determined by its dimension,
whereas the situation is more complicated over a general field such as
$\R$, $\Q_p$, or $\Q$.

% segment-id: o014.aljabr2.chapter7.galois-descent-h1.p013
Conditions imposed in the classification problems above, such as
nondegeneracy or associativity of multiplication, can usually be checked
after any field extension. The problem therefore splits into two parts:
% segment-id: o014.aljabr2.chapter7.galois-descent-h1.l004
\begin{enumerate}
	\item Over a sufficiently large field $L$, classify data of the form
	\eqref{eqn:tensor-data} satisfying the desired conditions; a typical choice
	is to take $L$ to be a separably closed field.
	\item Choose data $\mathbf{t}_0$ defined over $K$, take a sufficiently large
	Galois extension $L|K$, and study the data $\mathbf{t}$ satisfying
	$\mathbf{t}_L\simeq\mathbf{t}_{0,L}$.
\end{enumerate}
The first part may be called geometric, while the second is arithmetic. This
section focuses on the arithmetic part, which motivates the following
definition.

% segment-id: o014.aljabr2.chapter7.galois-descent-h1.d003
\begin{definition}
	Let $\mathbf{t}_0$ be data of the form \eqref{eqn:tensor-data} defined over
	the field $K$, and let $L|K$ be a Galois extension. Set
	% segment-id: o014.aljabr2.chapter7.galois-descent-h1.q012
	\[
	\mathcal{T}(\mathbf{t}_0):=
	\left\{
	\begin{array}{r|l}
		\mathbf{t}:\text{ data of the form \eqref{eqn:tensor-data}}
		& \text{there exists an isomorphism }h:\mathbf{t}_{0,L}\rightiso\mathbf{t}_L \\
		\text{defined over the field $K$}&
	\end{array}
	\right\}.
	\]
	Also write
	$\mathscr{T}(\mathbf{t}_0):=\mathcal{T}(\mathbf{t}_0)/\simeq$.
\end{definition}

% segment-id: o014.aljabr2.chapter7.galois-descent-h1.p014
This definition requires only the existence of an isomorphism $h$, without
specifying a choice; any two choices differ by right multiplication by an
element of $G(\mathbf{t}_{0,L})$. Write
% segment-id: o014.aljabr2.chapter7.galois-descent-h1.q013
\[
\Gamma:=\Gal(L|K),\qquad\text{equipped with the Krull topology}.
\]

% segment-id: o014.aljabr2.chapter7.galois-descent-h1.p015
If $\sigma\in\Gamma$ and
$h:\mathbf{t}_{0,L}\rightiso\mathbf{t}_L$ (that is,
$h(\mathbf{t}_{0,L})=\mathbf{t}_L$), then
% segment-id: o014.aljabr2.chapter7.galois-descent-h1.q014
\[
{}^\sigma h(\mathbf{t}_{0,L})
={}^\sigma h(\sigma\mathbf{t}_{0,L})
=\sigma(h(\mathbf{t}_{0,L}))
=\sigma(\mathbf{t}_L)
=\mathbf{t}_L.
\]
Consequently, there is a unique map
$c:\Gamma\to G(\mathbf{t}_{0,L})$ such that
% segment-id: o014.aljabr2.chapter7.galois-descent-h1.q015
\[
{}^\sigma h=hc(\sigma),\qquad \sigma\in\Gal(L|K).
\]

% segment-id: o014.aljabr2.chapter7.galois-descent-h1.p016
From
${}^{\sigma\tau}h={}^\sigma(hc(\tau))={}^\sigma h\cdot
{}^\sigma c(\tau)=h\cdot c(\sigma)\cdot{}^\sigma c(\tau)$, we obtain
% segment-id: o014.aljabr2.chapter7.galois-descent-h1.q016
\[
c(\sigma\tau)=c(\sigma)\cdot{}^\sigma c(\tau),
\qquad \sigma,\tau\in\Gamma.
\]

% segment-id: o014.aljabr2.chapter7.galois-descent-h1.p017
Without loss of generality, suppose that $\mathbf{t}_0$ and $\mathbf{t}$ are
realized on spaces $V$ and $W$, respectively, so that
$h\in\Hom_L(V_L,W_L)$. Since the action of $\Gamma$ is smooth, there is an
open subgroup $\Gamma_0\subset\Gamma$ that fixes $h$; consequently, $c$
factors through $\Gamma/\Gamma_0$.

% segment-id: o014.aljabr2.chapter7.galois-descent-h1.p018
If the choice of $h$ is changed to $h'=ha$, where
$a\in G(\mathbf{t}_{0,L})$, then the corresponding $c'$ satisfies
% segment-id: o014.aljabr2.chapter7.galois-descent-h1.q017
\[
c'(\sigma)=(h')^{-1}\cdot{}^\sigma h'
=a^{-1}\cdot h^{-1}\cdot{}^\sigma h\cdot{}^\sigma a
=a^{-1}\cdot c(\sigma)\cdot{}^\sigma a.
\]

% segment-id: o014.aljabr2.chapter7.galois-descent-h1.p019
Comparing these observations one by one with
Definition~\ref{def:nonabelian-H} of nonabelian $\Hm^1$, and using
Remark~\ref{rem:pro-finite-coh-nonabelian} for the profinite version, we
obtain a map
% segment-id: o014.aljabr2.chapter7.galois-descent-h1.q018
\begin{equation}\label{eqn:Galois-H1-prep}
	\begin{aligned}
		\mathcal{T}(\mathbf{t}_0)&\longrightarrow
		\Hm^1\left(\Gamma,G(\mathbf{t}_{0,L})\right),\\
		\mathbf{t}&\longmapsto[c]:=\text{the equivalence class of }c.
	\end{aligned}
\end{equation}
This map is canonical and independent of all auxiliary choices.

% segment-id: o014.aljabr2.chapter7.galois-descent-h1.p020
Recall also that $\Hm^1\left(\Gamma,G(\mathbf{t}_{0,L})\right)$ is generally
not a group, but a pointed set whose base point is represented by the
constant map $\Gamma\to G(\mathbf{t}_{0,L})$ with value the identity
automorphism $1$. On the other hand, $\mathcal{T}(\mathbf{t}_0)$ has the
obvious base point $\mathbf{t}_0$, as does its quotient
$\mathscr{T}(\mathbf{t}_0)$. The map \eqref{eqn:Galois-H1-prep} clearly
preserves base points, as one sees by taking $h=\identity$.

% segment-id: o014.aljabr2.chapter7.galois-descent-h1.t001
\begin{theorem}\label{prop:Galois-H1}
	For data $\mathbf{t}_0$ and a Galois extension $L|K$ as above, with
	$\Gamma:=\Gal(L|K)$, the map \eqref{eqn:Galois-H1-prep} induces a
	base-point-preserving bijection
	% segment-id: o014.aljabr2.chapter7.galois-descent-h1.q019
	\[
	\mathscr{T}(\mathbf{t}_0)\;\xrightarrow{1:1}\;
	\Hm^1\left(\Gamma,G(\mathbf{t}_{0,L})\right).
	\]
\end{theorem}
% segment-id: o014.aljabr2.chapter7.galois-descent-h1.pf003
\begin{proof}
	First we show that \eqref{eqn:Galois-H1-prep} factors through
	$\mathscr{T}(\mathbf{t}_0)$. Let
	$f:\mathbf{t}\rightiso\mathbf{t}'$, so that
	$f_L:\mathbf{t}_L\rightiso\mathbf{t}'_L$. From
	$h:\mathbf{t}_{0,L}\rightiso\mathbf{t}_L$ we obtain
	$f_Lh:\mathbf{t}_{0,L}\rightiso\mathbf{t}'_L$. Since
	${}^\sigma f_L=f_L$, we have
	$(f_Lh)^{-1}\cdot{}^\sigma(f_Lh)=h^{-1}\cdot{}^\sigma h$, and hence the
	corresponding $c$ is unchanged.

	Next we show that
	$\mathscr{T}(\mathbf{t}_0)\to
	\Hm^1\left(\Gamma,G(\mathbf{t}_{0,L})\right)$ is injective. Suppose that
	data $\mathbf{t}$ and $\mathbf{t}'$ and isomorphisms
	% segment-id: o014.aljabr2.chapter7.galois-descent-h1.q020
	\[
	\mathbf{t}'_L\xleftarrow[\sim]{h'}\mathbf{t}_{0,L}
	\xrightarrow[\sim]{h}\mathbf{t}_L,
	\]
	are given, and that there is an $a\in G(\mathbf{t}_{0,L})$ satisfying
	% segment-id: o014.aljabr2.chapter7.galois-descent-h1.q021
	\[
	(h')^{-1}\cdot{}^\sigma h'
	=a^{-1}\cdot h^{-1}\cdot{}^\sigma h\cdot{}^\sigma a,
	\qquad\sigma\in\Gamma.
	\]
	Consequently,
	% segment-id: o014.aljabr2.chapter7.galois-descent-h1.q022
	\[
	ha(h')^{-1}
	={}^\sigma h\cdot{}^\sigma a\cdot{}^\sigma(h')^{-1},
	\qquad\sigma\in\Gamma.
	\]
	In other words,
	$ha(h')^{-1}:\mathbf{t}'_L\rightiso\mathbf{t}_L$ is invariant under the
	$\Gamma$-action and therefore descends to an isomorphism
	$\mathbf{t}'\rightiso\mathbf{t}$.

	Finally, we show that
	$\mathscr{T}(\mathbf{t}_0)\to
	\Hm^1\left(\Gamma,G(\mathbf{t}_{0,L})\right)$ is surjective. Take
	$c\in Z^1(\Gamma,G(\mathbf{t}_{0,L}))$. Since
	$c(\sigma)\in\GL(V_L)$, we may use it to define a new $\Gamma$-action,
	denoted by $\odot$, on $V_L$:
	% segment-id: o014.aljabr2.chapter7.galois-descent-h1.q023
	\[
	\sigma\odot v:=c(\sigma)(\sigma v),
	\qquad\sigma\in\Gamma,\quad v\in V_L.
	\]

	From $c(\identity)=1$ and
	% segment-id: o014.aljabr2.chapter7.galois-descent-h1.q024
	\[
	\sigma\odot(\tau\odot v)
	=c(\sigma)\sigma(c(\tau)(\tau v))
	=c(\sigma)\bigl({}^\sigma c(\tau)(\sigma\tau v)\bigr)
	=(\sigma\tau)\odot v,
	\]
	it follows that $\odot$ is indeed a $\Gamma$-action, while the smoothness
	of $c$ ensures that $\odot$ is still a smooth semilinear action. Define the
	$K$-vector space
	$W:=(V_L)^{\Gamma,\;\odot\text{-action}}$. By Galois descent, we obtain an
	isomorphism of $L$-vector spaces that preserves the semilinear
	$\Gamma$-actions,
	% segment-id: o014.aljabr2.chapter7.galois-descent-h1.q025
	\begin{gather*}
		h:(V_L,\;\odot\text{-action})\rightiso
		(W_L,\;\text{natural action}),\\
		w\in W\implies h(w)=1\otimes w.
	\end{gather*}

	On the other hand, $\Gamma$ also acts semilinearly through $\odot$ on each
	$T^{p,q}(V_L)$ (see the discussion following
	Definition~\ref{def:tL}), and because
	$c(\sigma)\in G(\mathbf{t}_{0,L})$, we have
	% segment-id: o014.aljabr2.chapter7.galois-descent-h1.q026
	\[
	\sigma\odot\mathbf{t}_{0,L}
	=c(\sigma)(\sigma\mathbf{t}_{0,L})
	=c(\sigma)\mathbf{t}_{0,L}
	=\mathbf{t}_{0,L}.
	\]
	Thus, relative to $\odot$, the datum $\mathbf{t}_{0,L}$ descends to a datum
	$\mathbf{t}$ on $W$, characterized by
	$h(\mathbf{t}_{0,L})=\mathbf{t}_L$.

	It follows that $\mathbf{t}\in\mathcal{T}(\mathbf{t}_0)$. The corresponding
	isomorphism $h$ satisfies
	$\sigma(h(v))=h(\sigma\odot v)=(hc(\sigma)\sigma)(v)$, that is, the identity
	$\sigma h=hc(\sigma)\sigma$. Rewriting this as
	$h^{-1}\cdot{}^\sigma h=c(\sigma)$ shows that $\mathbf{t}$ maps to $[c]$
	under \eqref{eqn:Galois-H1-prep}. This proves the assertion.
\end{proof}

% segment-id: o014.aljabr2.chapter7.galois-descent-h1.p021
The simplest and most direct application is the following. It follows from
the general fact that every finite-dimensional vector space has a basis, so
that dimension determines the isomorphism class of a finite-dimensional
vector space.

% segment-id: o014.aljabr2.chapter7.galois-descent-h1.t002
\begin{theorem}[Hilbert's Theorem 90: the $\GL$ version]\label{prop:Hilbert-90-GL}
	\index{Hilbert 90}
	Let $V$ be a finite-dimensional $K$-vector space and let $L|K$ be a Galois
	extension. Thus $\Gamma:=\Gal(L|K)$ acts smoothly on the group $\GL(V_L)$.
	Then
	% segment-id: o014.aljabr2.chapter7.galois-descent-h1.q027
	\[
	\Hm^1(\Gamma,\GL(V_L))=1.
	\]
\end{theorem}
% segment-id: o014.aljabr2.chapter7.galois-descent-h1.pf004
\begin{proof}
	Apply Theorem~\ref{prop:Galois-H1}, but take data $\mathbf{t}_0$ with
	$I=\emptyset$. In other words, the data specify only a finite-dimensional
	$K$-vector space $V$, with no tensors. In this case, the elements of
	$\mathscr{T}(\mathbf{t}_0)$ correspond bijectively to the isomorphism
	classes of $K$-vector spaces $W$ satisfying $W_L\simeq V_L$, equivalently
	$\dim W=\dim V$. Clearly, there is only one such isomorphism class.
\end{proof}

% segment-id: o014.aljabr2.chapter7.galois-descent-h1.p022
Similarly, over a field of characteristic $\neq2$, the general fact that
every symplectic space $(V,\omega)$ has a symplectic basis leads to the
following result.

% segment-id: o014.aljabr2.chapter7.galois-descent-h1.t003
\begin{theorem}
	Let $\mathrm{char}(K)\neq2$, let $V$ be a $K$-vector space, and let
	$\omega:V\times V\to K$ be a nondegenerate antisymmetric bilinear form
	(also called a \emph{symplectic form}, while the datum $(V,\omega)$ is
	called a \emph{symplectic space}). The corresponding symplectic group is
	defined by
	% segment-id: o014.aljabr2.chapter7.galois-descent-h1.q028
	\begin{align*}
		\Sp(V)&=\Sp(V,\omega)\\
		&:=\{g\in\GL(V):\forall v_1,v_2\in V,\;
		\omega(gv_1,gv_2)=\omega(v_1,v_2)\}.
	\end{align*}

	Let $L|K$ be a Galois extension. Define the group $\Sp(V_L)$ and the
	action of $\Gamma:=\Gal(L|K)$ on it similarly; this action is smooth. Then
	% segment-id: o014.aljabr2.chapter7.galois-descent-h1.q029
	\[
	\Hm^1(\Gamma,\Sp(V_L))=1.
	\]
\end{theorem}
% segment-id: o014.aljabr2.chapter7.galois-descent-h1.pf005
\begin{proof}
	The data $\mathbf{t}_0$ under consideration consist of the space $V$ and
	the tensor $\omega$ of type $(0,2)$. The nondegeneracy and antisymmetry of
	a bilinear form can be checked after any field extension. Therefore the
	elements of $\mathscr{T}(\mathbf{t}_0)$ correspond bijectively to the
	isomorphism classes of symplectic spaces $(W,\eta)$ over $K$ satisfying
	$(W_L,\eta_L)\simeq(V_L,\omega_L)$. By the structure theorem for
	symplectic spaces---namely, that their dimension determines their
	isomorphism class---such data $(W,\eta)$ form a single isomorphism class.
\end{proof}

% segment-id: o014.aljabr2.chapter7.galois-descent-h1.eg001
\begin{example}\label{eg:multiplicative-90-finite}
	Take $\dim V=1$ in Theorem~\ref{prop:Hilbert-90-GL}. Since
	$\GL(V_L)\simeq L^\times$ with the standard $\Gamma$-action, we obtain
	% segment-id: o014.aljabr2.chapter7.galois-descent-h1.q030
	\[
	\Hm^1(\Gamma,L^\times)=0.
	\]
	Since $L^\times$ is abelian, additive notation is used in this equation.
	If $L|K$ is finite, it can be written in terms of the Tate cohomology of
	Definition~\ref{def:TaHm} as $\TaHm^1(\Gamma,L^\times)=0$.

	If $L|K$ is further assumed to be cyclic, the periodicity stated in
	Example~\ref{eg:TaHm-periodic} gives
	% segment-id: o014.aljabr2.chapter7.galois-descent-h1.q031
	\[
	\TaHm^{-1}(\Gamma,L^\times)
	\simeq\TaHm^1(\Gamma,L^\times)=0.
	\]
	Thus we naturally recover the multiplicative version of Hilbert's Theorem
	90; see \cite[Theorem 9.6.9]{Li1}.
\end{example}

% segment-id: o014.aljabr2.chapter7.galois-descent-h1.p023
In general, directly computing $\Hm^1$ in order to classify mathematical
objects may be rather difficult. Its value lies in the new perspective that
it provides, as well as in the availability of cohomological tools such as
the long exact sequence (Theorem~\ref{prop:nonabelian-long}) for studying
classification problems.
